\chapter{Disaster Response and Preparedness}\label{ch:disaster}

\section{General Disaster Principles}

When responding to disasters or mass casualty incidents (MCIs), personal safety and situational awareness take priority over individual rescue efforts.

\subsection{Scene Safety at Disasters}

\begin{enumerate}
  \item Assess for ongoing dangers: fire, explosion risk, structural collapse, hazardous materials, violence, electrical hazards
  \item Never enter an unstable structure or dangerous environment without proper equipment and training
  \item Establish a safety perimeter to prevent additional casualties
  \item Follow official instructions from emergency management personnel
  \item Use appropriate personal protective equipment when available (gloves, masks, eye protection)
\end{enumerate}

\noindent\textbf{WARNING}: Secondary incidents are common. Be aware of potential aftershocks (earthquakes), building collapses (fires/explosions), or renewed violence in active threat situations.

\section{Triage in Mass Casualty Incidents}

Triage is the process of sorting patients by urgency of need when resources are limited. The goal is to maximize survival for the greatest number of people \cite{who_mtbs}.

\subsection{The START Triage Method}

Simple Triage And Rapid Treatment (START) categorizes victims into four groups:

\begin{description}
  \item[Immediate (Red)] Life-threatening injuries requiring immediate treatment; transport priority 1
  \item[Delayed (Yellow)] Serious but not immediately life-threatening; can wait hours for treatment; transport priority 2
  \item[Minor (Green)] "Walking wounded"; minor injuries; transport priority 3
  \item[Deceased/Expectant (Black)] Deceased or expectant (injuries so severe that survival unlikely with available resources); transport last or none
\end{description}

\subsubsection{START Assessment Steps}

\begin{enumerate}
  \item \textbf{Walk Test}: Have all able-bodied individuals walk toward you. Those who cannot walk go to Red category
  \item \textbf{Respiratory Check}: For non-walking patients, assess breathing. If not breathing, open airway. If still not breathing -> Black (expectant/dead). If breathing proceed
  \item \textbf{Perfusion Check}: Check capillary refill (press fingernail bed — should color return in <2 seconds). If >2 seconds Red (poor perfusion/shock)
  \item \textbf{Mental Status Check}: Ask simple commands ("Shake your head," "Raise your hand"). If patient cannot obey Red (altered mental status). Otherwise Yellow (can follow commands, adequate breathing/perfusion)
\end{enumerate}

\noindent\textbf{NOTE}: Triage is dynamic. Patients categorized as Delayed may deteriorate and require reclassification to Immediate. Monitor all patients continuously.

\section{Specific Disaster Scenarios}

\subsection{Earthquakes}

\subsubsection{During an Earthquake}

\begin{itemize}
  \item If indoors: Drop, Cover, and Hold On — get under a sturdy desk/table, hold on until shaking stops. Stay away from windows, mirrors, heavy furniture that could fall
  \item If outdoors: Move away from buildings, power lines, and other hazards. Stay in open area until shaking stops
  \item If driving: Pull over to clear location (avoid bridges, overpasses, trees, signs). Stop safely and stay inside
\end{itemize}

\subsubsection{After the Shaking Stops}

\begin{itemize}
  \item Check for injuries and provide first aid
  \item Expect aftershocks — be prepared to "Drop, Cover, and Hold On" again
  \item Inspect for gas leaks (smell rotten eggs, hissing sound) — if detected, evacuate immediately and shut off main valve if safe
  \item Use telephones only for emergencies to keep lines clear for responders
  \item Listen to battery-powered radio for information
\end{itemize}

\subsection{Fire Emergencies}

\begin{itemize}
  \item Alert others and call emergency services immediately
  \item If small, contained fire and you have proper extinguisher type, attempt to extinguish only if escape route is clear behind you
  \item If unable to control fire, exit quickly using lowest position (smoke rises, cooler air stays near floor)
  \item Test doorknob with back of hand before opening — if hot, do NOT open
  \item If clothing catches fire: STOP, DROP, and ROLL (smother flames with coat if available)
  \item Once outside, stay outside — NEVER re-enter a burning building
\end{itemize}

\noindent\textbf{WARNING}: Smoke Inhalation. Victims of smoke inhalation may appear fine initially but develop respiratory distress hours later. Always seek medical evaluation after significant smoke exposure even if asymptomatic initially.

\subsection{Active Shooter or Violent Threat}

\begin{itemize}
  \item Run if there is an escape path — leave belongings behind, help others if possible, keep hands visible when exiting
  \item Hide if running is not possible — barricade door, silence phones, hide out of sight (do not make noise)
  \item Fight as a last resort — improvise weapons, act aggressively, commit to actions
\end{itemize}

\noindent\textbf{"RUN, HIDE, FIGHT"} is recommended by the U.S. Department of Homeland Security. Your priority is survival first, helping others second when it does not endanger yourself.

\subsection{Floods}

\begin{itemize}
  \item Never walk, swim, or drive through flowing floodwater -- 15 cm (6 inches) of moving water can knock a person down, and 30 cm (1 foot) can carry away a vehicle
  \item If a flood warning is issued, move to higher ground immediately; do not wait for water to rise
  \item Turn off gas and electricity at the mains if time and safety allow, and move valuables and important documents to upper floors
  \item Do not enter floodwater: it can contain sewage, chemicals, and debris, and may conceal submerged electrical wires
  \item After the water recedes, avoid contaminated floodwater in wounds; clean and dress any cuts and seek medical care if the wound was exposed to floodwater
  \item Discard food and water that came into contact with floodwater
\end{itemize}

\subsection{Cyclones (Hurricanes)}

Cyclones -- called hurricanes in the Atlantic and typhoons in the Pacific -- regularly strike India's coastal states (Odisha, Andhra Pradesh, Tamil Nadu, West Bengal, Gujarat, and the islands). The guidelines below apply to cyclones and hurricanes alike:

\begin{itemize}
  \item Prepare before the storm: secure outdoor objects, fill water containers, charge devices, and know the evacuation route and the nearest cyclone shelter
  \item Evacuate immediately when officials order it -- do not wait for the storm's worst; coastal districts are evacuated well in advance of landfall
  \item During the storm, stay inside in an interior room away from windows; use battery-powered radio for updates; the NDMA and state authorities broadcast cyclone warnings in local languages
  \item Beware the eye of the storm: winds may briefly calm, then resume violently from the opposite direction
  \item After the storm: stay clear of downed power lines and standing water near them; report injuries to emergency services and provide first aid as needed
\end{itemize}

\subsection{Heat Waves}

India experiences severe heat waves, mainly from April to June, when temperatures can exceed 45\textdegree C. Heat waves cause hundreds of deaths annually, largely among outdoor workers, the elderly, children, and people without cooling.

\begin{itemize}
  \item The National Disaster Management Authority (NDMA) publishes heat wave warnings and district-level action plans; follow official guidance and avoid going out between 12 noon and 3 p.m. during a heat wave
  \item Recognize heat exhaustion (headache, dizziness, nausea, heavy sweating, cramps) and heat stroke (hot dry skin, confusion, high body temperature, fainting) -- see Chapter~\ref{ch:environmental} for treatment
  \item If you must work outdoors, schedule frequent breaks in shade, drink water regularly without waiting for thirst, and use an umbrella or hat
  \item Check on elderly neighbours, the isolated, and those without fans or air conditioning; never leave children or pets in parked vehicles
  \item During a heat wave, community "cooling centres" are often opened in urban areas -- direct vulnerable people to them
\end{itemize}

\subsection{Landslides}

Landslides and mudflows occur in the Himalayan region, the Western Ghats, and the northeastern states, particularly after heavy or continuous rainfall.

\begin{itemize}
  \item During or after heavy rain, be alert in hilly areas for cracks in the ground, bulging soil, tilting trees or utility poles, and rumbling sounds -- these signal a possible landslide
  \item If you are indoors when a landslide starts, move to an upper floor or an interior room away from the hillside
  \item If outside, run uphill and away from the path of the landslide; do not stop to look back
  \item After a landslide, stay away from the area -- ground may shift further; report trapped or injured people to emergency services
\end{itemize}

\subsection{National Disaster Management Framework in India}

The \textbf{National Disaster Management Authority (NDMA)} coordinates disaster response across India through State and District Disaster Management Authorities and the National Disaster Response Force (NDRF) \cite{ndma}. The national disaster helpline is \textbf{1078}, and the unified emergency number \textbf{112} works in most states. Community volunteers and local bodies play a central role in first response.

\begin{itemize}
  \item Register your family's location with the local disaster management cell if you live in a hazard-prone area (cyclone, flood, earthquake, or landslide zone)
  \item Keep the NDMA helpline 1078 and the local emergency numbers (112/108) programmed in your phone
  \item During a disaster, follow the instructions of the NDRF, State Disaster Response Force, police, and civil authorities; their orders take priority over all other plans
  \item After a disaster, do not re-enter damaged buildings until authorities declare them safe
\end{itemize}

\subsection{Tornadoes}

\begin{itemize}
  \item Take shelter immediately when a tornado warning is issued
  \item Go to a basement or an interior room on the lowest floor, away from windows and corners
  \item In a vehicle or mobile home, leave immediately and find a sturdy building or low-lying area, lying flat and covering your head -- vehicles and mobile homes offer little protection
  \item Cover your head and neck with your arms or a heavy object (helmet, mattress, pillows)
  \item After the tornado, watch for downed power lines, gas leaks, broken glass, and unstable structures before moving
\end{itemize}

\subsection{Extreme Cold and Winter Storms}

\begin{itemize}
  \item Stay indoors and limit time outdoors during extreme cold warnings; if you must go out, dress in layers, cover exposed skin, and keep dry
  \item Know the signs of hypothermia (shivering, confusion, slurred speech) and frostbite (numbness, pale or waxy skin) -- see Chapter~\ref{ch:environmental}
  \item During a winter storm, keep a winter emergency kit in the vehicle: blankets, warm clothing, shovel, sand or cat litter, flashlight, water, and food
  \item If trapped in a vehicle: stay in the vehicle, run the engine and heater only in short intervals with a window cracked (check the exhaust is clear of snow to prevent carbon monoxide poisoning), and use the dome light sparingly
  \item Do not leave a running vehicle in a garage or fully enclosed space -- carbon monoxide is odorless and lethal
\end{itemize}

\subsection{Pandemics and Epidemics}

During infectious disease outbreaks, first aid principles remain the same, but additional precautions protect both responder and patient.

\begin{itemize}
  \item \textbf{Standard Precautions}: Treat all body fluids as potentially infectious. Wear gloves, use barrier devices for rescue breathing, and perform hand hygiene before and after care.
  \item \textbf{Source Control}: If the patient is able, have them wear a mask to reduce droplet spread. Provide masks to responders as well.
  \item \textbf{Respiratory Hygiene}: Cover coughs and sneezes with tissues or elbows. Dispose of tissues immediately and perform hand hygiene.
  \item \textbf{Surface Disinfection}: Clean and disinfect frequently touched surfaces. Use EPA-approved disinfectants or a dilute bleach solution (1 part bleach to 9 parts water) for contaminated areas.
  \item \textbf{PPE Selection}: In high-risk settings, add face shields, gowns, and N95/FFP2 respirators if available and trained to use.
  \item \textbf{Decontamination}: After providing care, remove PPE carefully (gloves first, then gown, then face protection, then hands) to avoid self-contamination. Perform hand hygiene immediately after.
  \item \textbf{Isolation until Evaluated}: If a person has fever and respiratory symptoms during an outbreak, minimize close contact, maintain distance when possible, and follow local health authority guidance.
  \item \textbf{Mental Health}: Outbreaks cause anxiety and stigma. Provide calm, factual information and psychological first aid to those affected.
\end{itemize}

\textbf{WARNING}: Do not provide first aid if you are at high risk for severe disease (elderly, immunocompromised, pregnant) and trained responders are unavailable -- call emergency services and guide the patient or bystanders through remote instructions if safe to do so.

\subsection{Loss of Power and Water}

\begin{itemize}
  \item Keep a battery-powered or hand-crank radio to receive official information
  \item Store at least 3 liters of water per person per day; if supplies run low, treat water from natural sources by boiling
  \item Keep refrigerators and freezers closed to preserve food; use coolers with ice for critical items
  \item Never use generators, camp stoves, or charcoal grills indoors -- they produce deadly carbon monoxide
  \item Keep a flashlight and spare batteries within reach, never use candles unattended
  \item Check on elderly or disabled neighbors who may be unable to prepare or call for help
\end{itemize}

\section{First Aid Kit Essentials}

A well-stocked first aid kit is critical preparation for any emergency \cite{cdc,arc_2020}. Contents should be stored in a waterproof, portable container and kept accessible.

A detailed checklist is provided in Appendix~\ref{ch:firstaid_kit}.

Regularly check expiration dates and replace used or expired items. Kits should be customized for the environment (home, workplace, vehicle, outdoor activities).

\section{Emergency Preparedness Checklist}

Every household should have an emergency plan and supplies ready:

\begin{itemize}
  \item Develop a family meeting plan (where to meet, how to contact if separated)
  \item Designate an out-of-area contact person (easier to reach long-distance during disasters)
  \item Assemble emergency kits (72-hour minimum supplies including water, food, medications, flashlight, batteries, radio, documents)
  \item Learn basic first aid and CPR certification (valid for 2 years; renew annually)
  \item Know local emergency procedures and evacuation routes
  \item Maintain updated copies of important documents (IDs, insurance policies, medical records) in waterproof containers
  \item Discuss special needs for infants, elderly, pets, and family members with disabilities
\end{itemize}

\section{Review Questions}

\begin{enumerate}[leftmargin=*]
\item What does the START triage method stand for, and what are its four outcome categories?
\item During a mass casualty incident, who gets treatment first: the person shouting loudly or the quiet one who is not responding?
\item Why do you remain in your car during an earthquake, and what should you do if you are indoors?
\item What is the correct order of actions during a fire?
\item Why is an out-of-area contact person recommended in a family emergency plan?
\item What minimum supplies should a 72-hour emergency kit contain?
\end{enumerate}

\index{mass casualty incident (MCI)}
\index{triage}
\index{START method}
\index{disaster response}
\index{earthquake safety}
\index{fire safety}
\index{active shooter response}
\index{emergency preparedness}
\index{first aid kit}
\index{supplies planning}
\index{family emergency plan}
\index{flood safety}
\index{hurricane safety}
\index{tornado safety}
\index{winter storm}
\index{power outage}
\index{carbon monoxide safety}
